\section*{Chapitre \ref{ch:s1:the-asset-managers-strategies} --- Les stratégies des sociétés de gestion}

\begin{solution}{exo:s1:the-asset-managers-strategies:1}
Soit $1/\sqrt{2} = 0.707$. Le tableau montre $4.84/2.96 = 1.64$ de 25 à 50 titres et $1.64/0.81 = 2.02$ de 100 à 200~: plus vite que la racine carrée,
car les titres ajoutés pèsent de moins en moins dans l'indice et le poids laissé de côté diminue plus vite que ne croît le nombre de titres.
\end{solution}

\begin{solution}{exo:s1:the-asset-managers-strategies:2}
$0.064 \times 10 = 0.64$ points de base par an.
\end{solution}

\begin{solution}{exo:s1:the-asset-managers-strategies:3}
La tracking error mesure l'écart des poids par rapport à l'indice, écart qui croît avec l'inclinaison~; le rendement dépend de ce que rapporte le facteur incliné, que le
produit ne contrôle pas.
\end{solution}

\begin{solution}{exo:s1:the-asset-managers-strategies:4}
$3.46~\%/\sqrt{252} = 0.218~\%$, soit 21.8 points de base par jour. Sur une transition d'un jour, toute la liste d'ordres est exposée pendant un jour à l'écart entre les
anciens et nouveaux portefeuilles, dont c'est l'écart-type quotidien.
\end{solution}

\begin{solution}{exo:s1:the-asset-managers-strategies:5}
Entre deux rééquilibrages, les poids détenus dérivent avec les rendements alors que les poids optimaux changent, des valeurs cotées entrent dans l'indice et en sortent, et le chiffre
ex ante est une moyenne, sur les dates de rééquilibrage, des prévisions faites avec les expositions du jour.
\end{solution}

\begin{solution}{exo:s1:the-asset-managers-strategies:6}
$0.03 \times 0.52 = 1.56~\%$ par an. Il a rapporté 0.07~\% parce que le ratio book-to-price, mesuré avec le cours du jour, se charge négativement sur le momentum introduit dans les données, dont
la prime compense la prime de valeur (chapitre~6).
\end{solution}

\begin{solution}{exo:s1:the-asset-managers-strategies:7}
Le coût plus deux écarts-types vaut 71.2, 70.0, 72.8, 80.3 et 98.3 points de base pour 1, 2, 3, 5 et 10 jours~: le planificateur choisit deux jours. Une
aversion plus forte au risque raccourcit la transition.
\end{solution}

\begin{solution}{exo:s1:the-asset-managers-strategies:8}
Un échantillon des 50 plus grands titres est un pari sur les grandes capitalisations contre le reste~; avec une tracking error voisine de 3~\%, dix ans à 2~\% par an restent à environ deux
erreurs types de zéro, et l'échantillon synthétique de 50 titres a fait de même sans aucune compétence. L'objectif de l'échantillonnage est une faible tracking error à faible coût, et non
un rendement excédentaire.
\end{solution}

\begin{solution}{pb:s1:the-asset-managers-strategies:1}
\begin{enumerate}
\item Détenir chaque composante au poids de l'indice~; détenir un sous-ensemble pondéré pour reproduire le risque de l'indice.
\item La tracking error est inévitable à cause des frictions, et les \omterm{def:s1:index-rebalancing:fund}{fonds indiciels} battent en moyenne les fonds actifs après frais.
\item Les poids pondérés par la capitalisation dérivent seuls avec les cours~; seules les entrées et les sorties nécessitent des transactions (6.4~\% par an).
\item Les expositions réelles du marché (bêta, secteurs, styles à chaque date) et les risques spécifiques~; un programme quadratique sur les $n$ plus grands titres minimisant la
tracking error ex ante, sans vente à découvert.
\item 4.84, 2.96, 1.64, 0.81 et 0.32~\% réalisés~; 4.70, 2.81, 1.52, 0.71 et 0.26~\% ex ante, pour 25 à 400 titres.
\item Leurs poids doivent être réoptimisés à mesure que l'indice évolue, et chaque réoptimisation déplace davantage de poids.
\item Rien~: un pari sur les grandes capitalisations dans la limite de sa tracking error.
\item Les fonds passifs profitent d'un rééquilibrage moins rigide~; les \omterm{def:s1:index-rebalancing:fund}{fonds indiciels} optimisés rééquilibrent plus tôt et plus patiemment, et paient moins.
\item Une inclinaison fondée sur des règles vers un style~; expositions 0.26, 0.52 et 1.03, tracking errors 1.57, 3.08 et 6.00~\%, rendements 0.06--0.07~\%.
\item La valeur mesurée avec le cours du jour revient à une position courte sur le momentum dans ce marché.
\item Prendre pour tracking error $\omega = \text{IR}/(2\lambda)$, avec IR le ratio d'information et $\lambda$ l'aversion à la variance active.
\item Déplacer les positions d'un portefeuille à l'autre au moindre implementation shortfall, en équilibrant l'impact de marché et le risque de la part non exécutée.
\item \textbf{Résultat nommé.} La \omterm{def:s1:the-asset-managers-strategies:replication}{réplication par échantillonnage} suit l'indice synthétique à 4.84~\%, 2.96~\%, 1.64~\%, 0.81~\% et 0.32~\% par an avec 25, 50, 100, 200 et 400
titres~; la transition de 59~\% d'un fonds de 2~milliards de dollars coûte 27.6 points de base en un jour ($\pm$21.8) et 18.5 en trois jours ($\pm$27.2), avec le coût plus un
écart-type minimal à deux ou trois jours (45.6).
\item L'impact de marché borne la vitesse~; le risque actif borne la patience.
\item Les titres communs aux deux portefeuilles, ou négociés contre d'autres flux, changent de propriétaire sans impact de marché.
\item Des erreurs ex ante plus grandes et un écart plus important avec la tracking error réalisée.
\item L'indice tel qu'il était, les dates d'annonce, les dividendes, les flux de trésorerie et les coûts.
\item Réplication de l'indice par échantillonnage.
\item Le chapitre~10 traitait les changements d'indice que les \omterm{def:s1:index-rebalancing:fund}{fonds indiciels} doivent suivre~; ici, le côté des fonds.
\item Garder un portefeuille proche d'une cible à moindre coût, et le déplacer sans perdre d'argent en chemin.
\end{enumerate}
\end{solution}

\begin{solution}{iq:s1:the-asset-managers-strategies:1}
Les coûts et frais, le frottement de trésorerie dû aux flux et aux dividendes, les transactions à des cours différents des cours de clôture de l'indice, l'échantillonnage au lieu de la \omterm{def:s1:the-asset-managers-strategies:replication}{réplication complète}, et
le choix du moment des transactions autour des changements d'indice.
\end{solution}

\begin{solution}{iq:s1:the-asset-managers-strategies:2}
Détenir intégralement les plus grands titres, puis choisir parmi les autres au moyen d'un optimiseur qui reproduit les expositions sectorielles, par pays et par style au moindre tracking error,
en respectant des limites de liquidité~; ou stratifier par secteur et par taille, puis retenir des titres représentatifs dans chaque cellule.
\end{solution}

\begin{solution}{iq:s1:the-asset-managers-strategies:3}
Construire la liste d'ordres, croiser d'abord les titres communs et les flux internes, couvrir l'exposition intermédiaire par des contrats à terme si la transition est longue, programmer le reste pour
équilibrer l'impact de marché et le risque actif, et rendre compte de l'implementation shortfall par rapport aux cours de décision.
\end{solution}

\begin{solution}{iq:s1:the-asset-managers-strategies:4}
L'annonce du fournisseur avec les ajouts, les suppressions et les changements de poids, les opérations sur titres, la date et l'heure d'effet, les tailles de transaction attendues par rapport à
la liquidité, et les contrôles vérifiant que les positions résultantes correspondent au nouvel indice à la tolérance près.
\end{solution}

\begin{solution}{iq:s1:the-asset-managers-strategies:5}
Son exposition au facteur visé et à ceux qu'il capte au passage (l'exposition au momentum d'un produit de valeur), la tracking error par rapport à l'indice parent, l'encombrement (crowding) dans
le facteur, et la liquidité~; il y a lieu de s'inquiéter lorsque les expositions non voulues augmentent ou que la corrélation du facteur avec d'autres se modifie.
\end{solution}

\begin{solution}{iq:s1:the-asset-managers-strategies:6}
$\sum_{d=0}^{N-1}(1 - d/N)^2 \approx N\int_0^1(1 - x)^2dx = N/3$, donc le risque vaut environ $\sigma_a\sqrt{N/3}$. En minimisant $c_0N^{-1/2} + \lambda\sigma_a
\sqrt{N/3}$ on obtient $-\tfrac12 c_0 N^{-3/2} + \tfrac{\lambda\sigma_a}{2\sqrt3}N^{-1/2} = 0$, donc $N = \sqrt3\,c_0/(\lambda\sigma_a)$.
\end{solution}
